\scanpage{6}
\diagram{\(\begin{array}{ll}
 E & \mathrm{Symm}_k(E^{\deleted{\unsure{\vee}}})^\vee\\
 P^r & \Phi^k(P^r).
\end{array}\)}
\margin{$\underbrace{k-1\quad k-1\qquad k-1}_{r+1\ \text{fois}}$}

\[
\deleted{\unsure{g}}\delta_{r,k}\in\Phi^{(k-1)(r+1)}\bigl(\Phi^k(P^r)\bigr)
\]
\margin{$(k-1)(r+1)$}
(discriminant) défini au signe près

Section de $\underline O_{\Phi^k(P^r)}((k-1)(r+1))$

Considérons une section \inserted{$\varphi$} de $\Phi^k(P^r)$ sur $S$, elle
définit donc une \deleted{\unsure{hom|homom.} $\Phi^k(P^r)$ section du}
faisceau inversible $\check M$ sur $S$,
% slanted
\margin{quotient de $\mathrm{Symm}_k(E)^\vee$ \unsure{sol.|i.e.} dual \unsure{ou|du} sous faisceau \unsure{qu'on prend} $M$ de $\mathrm{Symm}_k(E)$}
image inverse
de $\underline O_{\Phi^k(P^r)}(1)$] et \deleted{si l'image la section
est inversible par la} \inserted{$\delta_{r,k}$ a comme}
image réciproque une section $\underset{\vee}{\delta}$ de
$\check M^{(k-1)(r+1)}$. Dire la section $\varphi$ ne rencontre
pas $\Phi^k(P^r)_{\delta_{r,k}}$, i.e. $P^r_\varphi$ est simple sur $S$,
sss $\delta$ est partout non nulle. Notons que
d'ailleurs, si $J=J_\varphi$ est l'idéal défini par $\varphi$
sur $P^r$, alors
\[
J\deleted{\unsure{_\varphi}}\simeq g^*(M^{\deleted{\unsure{\vee}}})(-k)
\]
\margin{\deleted{\unsure{identifier|identifions} \unsure{ci-contre} $M$ -- $\check M$}}
\[
\begin{aligned}
[\,\deleted{O(-k)\otimes g^*(E)\to}\quad &\underline O(-k)\otimes g^*(\mathrm{Symm}_k(E))\longrightarrow\underline O_{P^r}\\
\text{d'où}\quad &\underline O(-k)\otimes g^*(M)\longrightarrow\underline O_{P^r}\quad -\,)
\end{aligned}
\]
\margin{\unsure{pour le} choix d'un}
% archive sheet number: 5

